% ═══════════════════════════════════════════════════════════════════════
% NTUA Greek Cover Pages — layout matched to peer Artemis template:
%   page 1 logo cover · page 2 blank · page 3 approval · page 4 copyright
% Included inside greek font / otherlanguage* from main_ntua.tex.
% ═══════════════════════════════════════════════════════════════════════

% ─────────────────────────────────────────────────────────────────────
% PAGE 1  —  Πρώτο Εξώφυλλο (με λογότυπο)
% ─────────────────────────────────────────────────────────────────────
\thispagestyle{empty}

\vspace*{-0.55cm}

\begin{center}
  \includegraphics[width=2.55cm]{ntua_logo.png}\\[0.45em]
  {\large\bfseries ΕΘΝΙΚΟ ΜΕΤΣΟΒΙΟ ΠΟΛΥΤΕΧΝΕΙΟ}\\[0.28em]
  {\normalsize ΣΧΟΛΗ ΗΛΕΚΤΡΟΛΟΓΩΝ ΜΗΧΑΝΙΚΩΝ ΚΑΙ ΜΗΧΑΝΙΚΩΝ ΥΠΟΛΟΓΙΣΤΩΝ}\\[0.18em]
  {\normalsize Τομέας Τεχνολογίας Πληροφορικής και Υπολογιστών}
\end{center}

\vspace{1.7cm}

\begin{center}
  {\large\bfseries Ανίχνευση Εισβολών Βασισμένη σε Τεχνητή Νοημοσύνη\\[0.28em]
         για Περιβάλλοντα Διαδικτύου Ιατρικών Πραγμάτων}\\[0.45em]
  {\normalsize\itshape AI Based Intrusion Detection for IoMT Environments}
\end{center}

\vspace{0.95cm}

\begin{center}
  {\large ΔΙΠΛΩΜΑΤΙΚΗ ΕΡΓΑΣΙΑ}
\end{center}

\vspace{0.55cm}

\begin{center}
  {\large\bfseries Παντελής Πετράκης}
\end{center}

\vfill

\noindent
\begin{tabular}{@{}ll@{}}
  \textbf{Επιβλέπων:} & Εμμανουήλ Βαρβαρίγος \\
                      & Καθηγητής Ε.Μ.Π.
\end{tabular}

\vspace{1.35cm}

\noindent Αθήνα, Σεπτέμβριος 2026

\clearpage

% ─────────────────────────────────────────────────────────────────────
% PAGE 2  —  Κενή σελίδα
% ─────────────────────────────────────────────────────────────────────
\thispagestyle{empty}
\null
\clearpage

% ─────────────────────────────────────────────────────────────────────
% PAGE 3  —  Δεύτερο Εξώφυλλο / Έγκριση
% ─────────────────────────────────────────────────────────────────────
\thispagestyle{empty}

\vspace*{-0.55cm}

\begin{center}
  {\large\bfseries ΕΘΝΙΚΟ ΜΕΤΣΟΒΙΟ ΠΟΛΥΤΕΧΝΕΙΟ}\\[0.28em]
  {\normalsize ΣΧΟΛΗ ΗΛΕΚΤΡΟΛΟΓΩΝ ΜΗΧΑΝΙΚΩΝ ΚΑΙ ΜΗΧΑΝΙΚΩΝ ΥΠΟΛΟΓΙΣΤΩΝ}\\[0.18em]
  {\normalsize Τομέας Τεχνολογίας Πληροφορικής και Υπολογιστών}
\end{center}

\vspace{1.5cm}

\begin{center}
  {\large\bfseries Ανίχνευση Εισβολών Βασισμένη σε Τεχνητή Νοημοσύνη\\[0.28em]
         για Περιβάλλοντα Διαδικτύου Ιατρικών Πραγμάτων}\\[0.45em]
  {\normalsize\itshape AI Based Intrusion Detection for IoMT Environments}
\end{center}

\vspace{0.85cm}

\begin{center}
  {\large ΔΙΠΛΩΜΑΤΙΚΗ ΕΡΓΑΣΙΑ}\\[0.55em]
  {\large\bfseries Παντελής Πετράκης}
\end{center}

\vspace{1.55cm}

\noindent
\begin{tabular}{@{}ll@{}}
  \textbf{Επιβλέπων:} & Εμμανουήλ Βαρβαρίγος \\
                      & Καθηγητής Ε.Μ.Π.
\end{tabular}

\vspace{1.15cm}

\noindent
Η διπλωματική εργασία εκπονήθηκε στη Facultat d'Informàtica de Barcelona (FIB)
του Universitat Politècnica de Catalunya (UPC), στο πλαίσιο του προγράμματος
Erasmus+, υπό την επίβλεψη του Καθηγητή Josep Pegueroles.

\vspace{1.15cm}

\begin{center}
  Εγκρίθηκε από την τριμελή επιτροπή την 21η Σεπτεμβρίου 2026.
\end{center}

\vfill

% Γραμμές υπογραφής + πλήρη ονόματα (3 στήλες, χωρίς wrap τίτλων)
\noindent
\begin{tabular*}{\linewidth}{@{}c@{\extracolsep{\fill}}c@{\extracolsep{\fill}}c@{}}
  \rule{4.2cm}{0.4pt} & \rule{4.2cm}{0.4pt} & \rule{4.2cm}{0.4pt} \\[0.45em]
  Εμμανουήλ Βαρβαρίγος & Θεοδώρα Βαρβαρίγου & Κωνσταντίνος Τσερπές \\
  {\footnotesize Καθηγητής Ε.Μ.Π.} &
  {\footnotesize Καθηγήτρια Ε.Μ.Π.} &
  {\footnotesize Επίκουρος Καθηγητής Ε.Μ.Π.}
\end{tabular*}

\vspace{0.85cm}

\begin{center}
  Αθήνα, Σεπτέμβριος 2026
\end{center}

\clearpage

% ─────────────────────────────────────────────────────────────────────
% PAGE 4  —  Copyright / Αποποίηση (δομή όπως πριν + υποχρεωτικό κείμενο Οδηγού)
% ─────────────────────────────────────────────────────────────────────
\thispagestyle{empty}

\vspace*{0.12\textheight}

\noindent\rule{\linewidth}{0.4pt}

\vspace{0.35em}

\noindent{\bfseries Pantelis Petrakis}\\[0.15em]
{\small
Διπλωματούχος Ηλεκτρολόγος Μηχανικός και Μηχανικός Ηλεκτρονικών
Υπολογιστών, Ε.Μ.Π.\\
Diploma in Electrical and Computer Engineering, NTUA
}

\vspace{0.55cm}

{\small
\noindent
Copyright \textcopyright\ Παντελής Πετράκης, 2026.\\
Με επιφύλαξη παντός δικαιώματος. All rights reserved.

\vspace{0.4cm}

\noindent
Απαγορεύεται η αντιγραφή, αποθήκευση και διανομή της παρούσας εργασίας,
εξ ολοκλήρου ή τμήματος αυτής, για εμπορικό σκοπό.
Επιτρέπεται η ανατύπωση, αποθήκευση και διανομή για σκοπό μη κερδοσκοπικό,
εκπαιδευτικής ή ερευνητικής φύσης, υπό την προϋπόθεση να αναφέρεται η πηγή
προέλευσης και να διατηρείται το παρόν μήνυμα.
Ερωτήματα που αφορούν τη χρήση της εργασίας για κερδοσκοπικό σκοπό πρέπει
να απευθύνονται προς τον συγγραφέα.

\vspace{0.4cm}

\noindent
Οι απόψεις και τα συμπεράσματα που περιέχονται σε αυτό το έγγραφο εκφράζουν
τον συγγραφέα και δεν πρέπει να ερμηνευθεί ότι αντιπροσωπεύουν τις επίσημες
θέσεις του Εθνικού Μετσόβιου Πολυτεχνείου.
}

\vspace{0.55cm}

\noindent{\itshape\bfseries Pantelis Petrakis}\\[0.1em]
{\small\itshape Diploma in Electrical and Computer Engineering, NTUA.}

\vspace{0.35cm}

{\small\itshape
\noindent
Copyright \textcopyright\ Pantelis Petrakis, 2026. All Rights Reserved.

\vspace{0.35cm}

\noindent
Copying, storage and distribution of this thesis, in whole or in part, for
commercial purposes is prohibited.
Reproduction, storage and distribution for non-profit, educational or research
purposes is permitted, provided that the source is mentioned and this message
is retained.
Enquiries concerning the use of this work for profit-making purposes must be
addressed to the author.

\vspace{0.35cm}

\noindent
The views and conclusions contained in this document are those of the author and
should not be interpreted as representing the official positions of the National
Technical University of Athens.
}

\clearpage
